% Lösungen Einheit 3
\einheitenkopf{Einheit 3 von 4 · Normalform und p-q-Formel}
\zweigzeile{Hier lernst du, eine quadratische Gleichung in die Normalform zu bringen und mit der p-q-Formel zu lösen · neu in diesem Jahr (an der Oberschule je nach Buch Klasse 9 oder 10) · P10 oft · baut auf: Wurzelziehen und Lösbarkeit (Einheit 1), Satz vom Nullprodukt (Einheit 2), Terme ordnen, binomische Formeln}
\erg{25}{a) ohne $x$-Glied, Wurzelziehen \quad b) Produkt $= 0$, Nullprodukt \quad c) Klammer$^2 =$ Zahl, Rückwärtsrechnen \quad d) Normalform, Formel \quad e) ohne $x$-Glied, Wurzelziehen}
\erg{26}{a) $p = 5$, $q = 4$ \quad b) $p = -3$, $q = 2$ \quad c) $p = 1$, $q = -12$ \quad d) $p = -11$, $q = 0$ \quad e) $p = 0$, $q = -2$}
\erg{27}{a) teilen durch $2$ \quad b) nicht teilen \quad c) teilen durch $-1$ \quad d) teilen durch $\frac{1}{2}$ (also mal $2$) \quad e) teilen durch $5$}
\erg{28}{a) $-3 \pm 2$; $x_1 = -1$, $x_2 = -5$ \quad b) $-5 \pm 2$; $x_1 = -3$, $x_2 = -7$ \quad c) $-6 \pm 1$; $x_1 = -5$, $x_2 = -7$ \quad d) $-7 \pm 3$; $x_1 = -4$, $x_2 = -10$ \quad e) $2 \pm 3$; $x_1 = 5$, $x_2 = -1$}
\erg{29}{a) $x^2 - 2x - 24 = 0$; $1 \pm 5$; $x_1 = 6$, $x_2 = -4$ \quad b) $: (-2)$: $x^2 - 6x + 5 = 0$; $3 \pm 2$; $x_1 = 5$, $x_2 = 1$ \quad c) $2{,}5 \pm \sqrt{6{,}25 - 6} = 2{,}5 \pm 0{,}5$; $x_1 = 3$, $x_2 = 2$ \quad d) $5 \pm \sqrt{25 - 25} = 5$; eine Lösung: $x = 5$ \quad e) $\sqrt{1 - 7} = \sqrt{-6}$; keine Lösung \quad f) $D = 12{,}25 - 11 = 1{,}25$; positiv, also zwei Lösungen}
\erg{30}{a) $2 \pm \sqrt{3}$; $x_1 \approx 3{,}73$, $x_2 \approx 0{,}27$ \quad b) $49 - 35 - 14 = 0$, $0 = 0$ (wA) \quad c) $x^2 + 6x + 9 = 2x + 14$, $x^2 + 4x - 5 = 0$; $-2 \pm 3$; $x_1 = 1$, $x_2 = -5$}
\erg{31}{a) Wurzelziehen: $x^2 = 25$; $x_1 = 5$, $x_2 = -5$ \quad b) Nullprodukt: $x_1 = 6$, $x_2 = -1$ \quad c) Nullprodukt nach Ausklammern: $x_1 = 0$, $x_2 = 9$ \quad d) Rückwärtsrechnen: $x - 4 = 1$ oder $x - 4 = -1$; $x_1 = 5$, $x_2 = 3$ \quad e) Formel: $-0{,}5 \pm 2{,}5$; $x_1 = 2$, $x_2 = -3$ \quad f) $2x^2 + 4x - 4 = 0 \mid :2$; $x^2 + 2x - 2 = 0$; $x_{1,2} = -1 \pm \sqrt{1 + 2} = -1 \pm \sqrt{3}$; $x_1 \approx 0{,}73$, $x_2 \approx -2{,}73$}
\erg{32}{a) $x^2 - 2x - 8 = 0$; $1 \pm 3$; $x_1 = 4$, $x_2 = -2$; $S_1(4|13)$, $S_2(-2|1)$ \quad b) $x^2 - 6x + 9 - 2 = -x + 3$, $x^2 - 5x + 4 = 0$; $2{,}5 \pm 1{,}5$; $x_1 = 4$, $x_2 = 1$; $S_1(4|-1)$, $S_2(1|2)$}
\erg{33}{a) Fehler: Vorzeichen von $-\frac{p}{2}$ – bei $p = -14$ ist $-\frac{p}{2} = 7$, nicht $-7$; richtig: $7 \pm 2$; $x_1 = 9$, $x_2 = 5$ \quad b) $4 \pm 3$; $x_1 = 7$, $x_2 = 1$}
\erg{34}{a) die Formel gilt nur für die Normalform, in der $x^2$ die Vorzahl eins hat; ohne Teilen wären $p$ und $q$ falsch \quad b) $\left(\frac{p}{2}\right)^2$ ist nie negativ, und $-q$ ist positiv; der Wert unter der Wurzel ist also positiv, und es gibt zwei Lösungen}
\erg{35}{a) $-0{,}1x^2 + 0{,}6x + 1{,}6 = 0 \mid :(-0{,}1)$; $x^2 - 6x - 16 = 0$; $3 \pm 5$; $x_1 = 8$, $x_2 = -2$; der Ball kommt nach $8\,\text{m}$ auf, $-2$ passt nicht (vor dem Abwurf) \quad b) $-0{,}1x^2 + 0{,}6x + 1{,}6 = 2$; $x^2 - 6x + 4 = 0$; $3 \pm \sqrt{5}$; nach etwa $0{,}76\,\text{m}$ und $5{,}24\,\text{m}$}
